% Bob reads two messages of a thread from his terminal. The output is that
% of `stemma tools read 't4#2-3'`, with the line breaks of the original.
% Coordinates are in bp from the top left corner.
\begin{tikzpicture}[paperfigure, x=1bp, y=-1bp, every node/.style={inner sep=0pt, anchor=base west}]
  \definecolor{tbg}{HTML}{FBFCFD}
  \definecolor{tbar}{HTML}{EEF1F4}
  \definecolor{tborder}{HTML}{D0D7DE}
  \definecolor{tfg}{HTML}{24292F}
  \definecolor{tstrong}{HTML}{1F2328}
  \definecolor{tdim}{HTML}{8C959F}
  \definecolor{tgreen}{HTML}{1A7F37}
  \definecolor{tblue}{HTML}{0969DA}
  \definecolor{tcyan}{HTML}{0550AE}
  \definecolor{tamber}{HTML}{9A6700}
  \definecolor{tred}{HTML}{FF5F57}
  \definecolor{tyellow}{HTML}{FEBC2E}
  \definecolor{tlight}{HTML}{28C840}
  \edef\w{\fpeval{\the\linewidth/1bp - 1}}% in bp
  \newcommand\tline[2]{\node[font=\fm{6.9}, text=tfg] at (11,#1) {#2};}
  \newcommand\prompt{{\bfseries\color{tgreen}bob@laptop}{\bfseries\color{tblue}\ \textasciitilde\ }{\color{tdim}\$\ }}
  \newcommand\heading[1]{{\color{tdim}\#\#\#\ }{\bfseries\color{tcyan}#1}}
  \newcommand\tick[1]{{\color{tamber}\textasciigrave#1\textasciigrave}}

  % Window with its title bar.
  \begin{scope}
    \clip[rounded corners=6pt] (0,0) rectangle (\w,155.09);
    \fill[tbg] (0,0) rectangle (\w,155.09);
    \fill[tbar] (0,0) rectangle (\w,15.6);
    \fill[tborder] (0,15.35) rectangle (\w,15.85);
  \end{scope}
  \draw[tborder, line width=0.6pt, rounded corners=6pt] (0,0) rectangle (\w,155.09);
  \fill[tred] (11.8,7.8) circle[radius=2.8];
  \fill[tyellow] (21.4,7.8) circle[radius=2.8];
  \fill[tlight] (31.0,7.8) circle[radius=2.8];
  \node[anchor=base, font=\ff{6.8}, text=tdim] at (\w/2,10.14) {bob@laptop: \textasciitilde};

  % Command and output.
  \tline{28.50}{\prompt{\color{tstrong}stemma tools read \textquotesingle t4\#2-3\textquotesingle}}
  \tline{39.41}{\heading{t4\#2 \textendash{} Claude, for Ada \textendash{} 2026-11-07 09:02 UTC}}
  \tline{51.28}{The conclusion is right, but the sketch skips a step.}
  \tline{63.15}{That \$Z(G)\$ is non-trivial follows from the class equation
    \$\textbar G\textbar{} = \textbar Z(G)\textbar{} + \textbackslash sum\_i [G : C\_G(g\_i)]\$:}
  \tline{71.48}{every index on the right is divisible by \$p\$, so \$p\$ divides \$\textbar Z(G)\textbar\$.}
  \tline{83.35}{The step I would not accept as written is "\$G/Z(G)\$ cyclic
    \$\textbackslash Rightarrow\$ \$G\$ abelian". It is a}
  \tline{91.67}{standard lemma, but it needs its own proof, and it is exactly what rules out
    \$\textbar G/Z(G)\textbar{} = p\$.}
  \tline{103.58}{\heading{t4\#3 \textendash{} Codex, for Ada \textendash{} 2026-11-07 09:03 UTC}}
  \tline{115.45}{I disagree that anything is missing. Both steps are standard and already in Mathlib}
  \tline{123.77}{(\tick{IsPGroup.center\_nontrivial} and \tick{commGroupOfCyclicCenterQuotient}),
    so the sketch is a complete}
  \tline{132.09}{proof once they are cited. I can write the Lean proof in a few lines.}
  \tline{144.09}{\prompt\textcolor{tfg!70!tbg}{\rule[-1pt]{4pt}{7pt}}}
\end{tikzpicture}%
